\expandafter\def\csname HistoricalNoteText3\endcsname{Der analytischen Vollständigkeit wegen bemerke ich, dass in diesem Falle der Himmelskörper in
$\left\{\begin{array}{c}\AscendingNode\\\DescendingNode\end{array}\right\}$ ist, je nachdem $\sin(\odot-\alpha)$ und $\sin(\alpha-\AscendingNode)$
$\left\{\begin{array}{l}\text{einerlei}\\\text{entgegengesetzte}\end{array}\right\}$ Zeichen haben.}
\clearpage\SourceLeaf{92}%
\pdfbookmark[1]{Vorschriften zur Bestimmung heliocentrischer Orte}{vorschriften1802}
\begin{center}
{\small VORSCHRIFTEN, UM AUS DER GEOCENTRISCHEN LÄNGE UND BREITE EINES
HIMMELSKÖRPERS, DEM ORTE SEINES KNOTENS, DER NEIGUNG DER BAHN, DER LÄNGE
DER SONNE UND IHREM ABSTANDE VON DER ERDE ABZULEITEN:
DES HIMMELSKÖRPERS HELIOCENTRISCHE LÄNGE IN DER BAHN, WAHREN ABSTAND
VON DER SONNE UND WAHREN ABSTAND VON DER ERDE.\par}
\medskip\hrule\medskip
{\small Monatliche Correspondenz zur Beförderung der Erd- und Himmels-Kunde,\par
herausgegeben vom Freiherrn \textsc{von Zach}. Juni 1802.\par}
\medskip\hrule
\end{center}\medskip
\begin{center}\textit{Bedeutung der Zeichen:}\end{center}
\noindent\begin{minipage}[t]{.48\textwidth}\small
\begin{center}\textit{Gegeben;}\end{center}
\begin{tabular}{@{}lp{.78\linewidth}@{}}
$\AscendingNode$ & \raggedright Länge des aufsteigenden Knotens.\tabularnewline
$\odot$ & \raggedright Länge der Sonne.\tabularnewline
$\alpha$ & \raggedright Geocentrische Länge des Himmelskörpers.\tabularnewline
$\ObservedBeta$ & \raggedright Geocentrische Breite.\tabularnewline
$i$ & \raggedright Neigung der Bahn.\tabularnewline
$R$ & \raggedright Abstand der Sonne von der Erde.
\end{tabular}
\end{minipage}\hfill
\begin{minipage}[t]{.48\textwidth}\small
\begin{center}\textit{Gesucht:}\end{center}
\begin{tabular}{@{}lp{.78\linewidth}@{}}
$\upsilon$ & \raggedright heliocentr. Länge des Himmelskörpers in der Bahn.\tabularnewline
$r$ & \raggedright Wahrer Abstand von der Sonne.\tabularnewline
$\Delta$ & \raggedright Wahrer Abstand von der Erde.
\end{tabular}
\[
\left.\begin{array}{l}A\\B\\C\ \text{etc.}\end{array}\right\}\text{ Hülfswinkel.}
\]
\end{minipage}

\begin{center}I.\end{center}
\[
\begin{aligned}
1^{0}\quad\frac{\cos(\odot-\AscendingNode)\operatorname{tang}\ObservedBeta}{\sin(\odot-\alpha)}&=\operatorname{tang}A\\[5pt]
\frac{\sin A\operatorname{tang}(\odot-\AscendingNode)}{\sin(A+i)}&=\operatorname{tang}(\upsilon-\AscendingNode)
\end{aligned}
\]
\[
\begin{aligned}
2^{0}\quad\frac{\sin(\odot-\alpha)\operatorname{tang}i}{\cos(\odot-\AscendingNode)}&=\operatorname{tang}B\\[5pt]
\frac{\cos B\sin\ObservedBeta\operatorname{tang}(\odot-\AscendingNode)}{\sin(B+\ObservedBeta)\cos i}&=\operatorname{tang}(\upsilon-\AscendingNode)
\end{aligned}
\]
\[
\begin{aligned}
3^{0}\quad\frac{\sin(\odot-\AscendingNode)\operatorname{tang}\ObservedBeta}{\sin(\odot-\alpha)\operatorname{tang}i}&=\operatorname{tang}C\\[5pt]
\frac{\sin C\sin(\odot-\AscendingNode)}{\sin(C+\odot-\AscendingNode)\cos i}&=\operatorname{tang}(\upsilon-\AscendingNode)
\end{aligned}
\]
\[
\begin{aligned}
4^{0}\quad\frac{\cos(\odot-\AscendingNode)\operatorname{tang}\ObservedBeta}{\cos(\odot-\alpha)\operatorname{tang}i}&=\operatorname{tang}D\\[5pt]
\frac{\sin D\operatorname{tang}(\odot-\AscendingNode)\cos(\odot-\alpha)}{\sin(D+\odot-\alpha)\cos i}&=\operatorname{tang}(\upsilon-\AscendingNode)
\end{aligned}
\]
\textit{Anmerkung:} Da Winkel, die um $180^{0}$ verschieden sind, einerlei Tangenten
haben, so ist hier noch eine Vorschrift nöthig, wie die durch ihre Tangenten
bestimmten Winkel $A$, $B$, $C$ etc. und $\upsilon-\AscendingNode$ angesetzt werden müssen. Den
Winkel $\upsilon-\AscendingNode$ hat man allezeit zwischen 0 und $180^{0}$ anzunehmen, wenn $\ObservedBeta$
\SourceLeaf{93}%
positiv (nördlich) ist; ist hingegen die Breite südlich, so muss $\upsilon-\AscendingNode$ zwischen
$180^{0}$ und $360^{0}$, oder was einerlei ist, zwischen $-180^{0}$ und 0 fallen. Ist
$\ObservedBeta=0$, so ist der Himmelskörper in einem Knoten, und man wird nie zweifelhaft
sein, ob es $\AscendingNode$ oder $\DescendingNode$ ist\HistoricalNote{3}. Die Hülfswinkel, $A$, $B$, $C$, $D$ aber, so wie
die folgenden $E$, $F$ etc., kann man in dieser Hinsicht ganz nach Belieben ansetzen;
wobei es sich jedoch von selbst versteht, dass man auf die Zeichen $\pm$
gehörige Rücksicht nehme; ich habe sie in folgendem Beispiele immer zwischen
$-90^{0}$ und $+90^{0}$ genommen. — Logarithmen, deren zugehörige Grössen negativ
sind, habe ich durch ein beigeschriebenes n ausgezeichnet.

\begin{center}II.\end{center}
\[
\begin{aligned}
5^{0}\quad\frac{\operatorname{tang}\ObservedBeta}{\sin(\alpha-\AscendingNode)}&=\operatorname{tang}E\\[5pt]
\frac{\sin E\sin(\odot-\AscendingNode)}{\sin(i-E)\sin(\upsilon-\AscendingNode)}&=\frac{r}{R}
\end{aligned}
\]
\[
\begin{aligned}
6^{0}\quad\operatorname{tang}i\sin(\alpha-\AscendingNode)&=\operatorname{tang}F\\[5pt]
\frac{\cos F\sin(\odot-\AscendingNode)\sin\ObservedBeta}{\sin(F-\ObservedBeta)\sin(\upsilon-\AscendingNode)\cos i}&=\frac{r}{R}
\end{aligned}
\]
\[
\begin{aligned}
7^{0}\quad\cos i\operatorname{tang}(\upsilon-\AscendingNode)&=\operatorname{tang}G\\[5pt]
\frac{\cos G\sin(\odot-\alpha)}{\sin(\alpha-\AscendingNode-G)\cos(\upsilon-\AscendingNode)}&=\frac{r}{R}
\end{aligned}
\]
\[
\begin{aligned}
8^{0}\quad\frac{\operatorname{tang}(\alpha-\AscendingNode)}{\cos i}&=\operatorname{tang}H\\[5pt]
\frac{\sin H\sin(\odot-\alpha)}{\sin(H-(\upsilon-\AscendingNode))\sin(\alpha-\AscendingNode)}&=\frac{r}{R}
\end{aligned}
\]
\[
\begin{aligned}
9^{0}\quad\frac{\operatorname{tang}\ObservedBeta}{\sin i\cos(\alpha-\AscendingNode)}&=\operatorname{tang}I\\[5pt]
\frac{\sin I\cos(\odot-\AscendingNode)}{\sin(\upsilon-\AscendingNode-I)}&=\frac{r}{R}
\end{aligned}
\]
\[
\begin{aligned}
10^{0}\quad\sin i\cos(\alpha-\AscendingNode)\operatorname{tang}(\upsilon-\AscendingNode)&=\operatorname{tang}K\\[5pt]
\frac{\cos K\sin\ObservedBeta\cos(\odot-\AscendingNode)}{\sin(K-\ObservedBeta)\cos(\upsilon-\AscendingNode)}&=\frac{r}{R}
\end{aligned}
\]
\[
\begin{aligned}
11^{0}\quad\frac{\sin C\sin(\odot-\alpha)}{\cos(C+\odot-\alpha)\operatorname{tang}(\odot-\AscendingNode)\cos i}&=\operatorname{tang}L\\[5pt]
\frac{\sin L}{\sin(\upsilon-\AscendingNode-L)\cos(\odot-\AscendingNode)}&=\frac{r}{R}
\end{aligned}
\]
\[
\begin{aligned}
12^{0}\quad\frac{\sin D\cos(\odot-\AscendingNode)}{\cos(D+\odot-\AscendingNode)\cos i}&=\operatorname{tang}M\\[5pt]
\frac{\sin M}{\sin(\upsilon-\AscendingNode-M)\cos(\odot-\AscendingNode)}&=\frac{r}{R}
\end{aligned}
\]
\begin{center}III.\end{center}
\[
13^{0}\quad\frac{r\sin(\upsilon-\AscendingNode)\sin i}{\sin\ObservedBeta}=\Delta
\]
\[
\begin{aligned}
14^{0}\quad\frac{R\sin E\sin(\odot-\AscendingNode)\sin i}{\sin(i-E)\sin\ObservedBeta}&\\[5pt]
&=\frac{R\cos E\sin(\odot-\AscendingNode)\sin i}{\sin(i-E)\sin(\alpha-\AscendingNode)\cos\ObservedBeta}\\[5pt]
&=\Delta
\end{aligned}
\]
\[
\begin{aligned}
15^{0}\quad\frac{R\cos F\sin(\odot-\AscendingNode)\operatorname{tang}i}{\sin(F-\ObservedBeta)}&\\[5pt]
&=\frac{R\sin F\sin(\odot-\AscendingNode)\sin(\alpha-\AscendingNode)}{\sin(F-\ObservedBeta)}\\[5pt]
&=\Delta
\end{aligned}
\]
Und so lassen sich noch mehrere Ausdrücke für $\Delta$ aus der Verbindung
von $13^{0}$ mit allen Formeln II ableiten.

\SourceLeaf{94}%
\begin{center}\textit{Beispiel.}\end{center}
\[
\begin{array}{rlrrr}
\AscendingNode&=&80^{0}&59'&12''07\\
\odot&=&281&1&34.99\\
\alpha&=&53&23&2.46\\
i&=&10&37&9.55
\end{array}
\]
\[
\log\operatorname{tang}\ObservedBeta=8.7349698\,\mathrm n
\qquad\log R=9.9926158
\]
\[
\ObservedBeta=-3^{0}\ 6'\ 33''561\quad\text{negativ oder südlich}
\]
\[
\begin{array}{rlrrr}
\text{Folglich }\odot-\AscendingNode&=&200^{0}&2'&22''92\\
\odot-\alpha&=&227&38&32.53\\
\alpha-\AscendingNode&=&-27&36&9.61
\end{array}
\]
\par\medskip\noindent\begin{minipage}[t]{\textwidth}
\begin{center}$1^{0}$.\end{center}
\begin{center}\small
\begin{tabular}{@{}lcl@{}}
$\log\operatorname{tang}\ObservedBeta$ & $ .\ .\ .\ .\ . $ & $8.7349698\,\mathrm n$\\
$\log\cos(\odot-\AscendingNode)$ & $ .\ .\ . $ & $9.9728762\,\mathrm n$\\
$\text{Compl.}\ \log\sin(\odot-\alpha)$ & $ . $ & $0.1313827\,\mathrm n$\\\hline
$\log\operatorname{tang}A$ & $ .\ .\ .\ .\ . $ & $8.8392287\,\mathrm n$\\\hline
$\log\sin A$ & $ .\ .\ .\ .\ . $ & $8.8381955\,\mathrm n$\\
$\log\operatorname{tang}(\odot-\AscendingNode)$ & $ .\ .\ . $ & $9.5620014$\\
$\text{Compl.}\ \log\sin(A+i)$ & $ . $ & $0.9350608$\\\hline
$\log\operatorname{tang}(\upsilon-\AscendingNode)$ & $ .\ .\ . $ & $9.3352577\,\mathrm n$
\end{tabular}
\end{center}
\[
\begin{array}{llrrr}
\text{folglich}&A=-&3^{0}&57'&2''136\\
&A+i=&6&40&7.414\\[4pt]
\text{ferner}&\upsilon-\AscendingNode=-&12^{0}&12'&37''942\\
\text{also}&\upsilon=.&68&46&34.128\\\hline
\end{array}
\]

\end{minipage}\par
\par\medskip\noindent\begin{minipage}[t]{\textwidth}
\begin{center}$2^{0}$.\end{center}
\begin{center}\small
\begin{tabular}{@{}lcl@{}}
$\log\sin(\odot-\alpha)$ & $ .\ .\ .\ . $ & $9.8686173\,\mathrm n$\\
$\log\operatorname{tang}i$ & $ .\ .\ .\ .\ . $ & $9.2729872$\\
$\text{Compl.}\ \log\cos(\odot-\AscendingNode)$ & $ . $ & $0.0271238\,\mathrm n$\\\hline
$\log\operatorname{tang}B$ & $ .\ .\ .\ .\ . $ & $9.1687283$\\\hline
$\log\cos B$ & $ .\ .\ .\ .\ . $ & $9.9953277$\\
$\log\sin\ObservedBeta$ & $ .\ .\ .\ .\ .\ . $ & $8.7343300\,\mathrm n$\\
$\log\operatorname{tang}(\odot-\AscendingNode)$ & $ .\ .\ . $ & $9.5620014$\\
$\text{Compl.}\ \log\sin(B+\ObservedBeta)$ & $ . $ & $1.0360961$\\
$\text{Compl.}\ \log\cos i$ & $ .\ .\ . $ & $0.0075025$\\\hline
$\log\operatorname{tang}(\upsilon-\AscendingNode)$ & $ .\ .\ . $ & $9.3352577\,\mathrm n$
\end{tabular}
\end{center}
\[
\begin{array}{llrrr}
\text{folglich}&B=&8^{0}&23'&21''888\\
&B+\ObservedBeta=&5&16&48.327
\end{array}
\]
wie oben

\end{minipage}\par

\SourceLeaf{95}%

\par\medskip\noindent\begin{minipage}[t]{\textwidth}
\begin{center}$3^{0}$.\end{center}
\begin{center}\small
\begin{tabular}{@{}lcl@{}}
$\log\sin(\odot-\AscendingNode)$ & $.\ .\ .$ & $9.5348776\,\mathrm n$\\
$\log\operatorname{tang}\ObservedBeta$ & $.\ .\ .\ .\ .$ & $8.7349698\,\mathrm n$\\
$\text{Compl.}\ \log\sin(\odot-\alpha)$ & $.$ & $0.1313827\,\mathrm n$\\
$\text{Compl.}\ \log\operatorname{tang}i$ & $.\ .\ .$ & $0.7270128$\\
\hline
$\log\operatorname{tang}C$ & $.\ .\ .\ .\ .$ & $9.1282429\,\mathrm n$\\
\hline
$\log\sin C$ & $.\ .\ .\ .\ .$ & $9.1243583\,\mathrm n$\\
$\log\sin(\odot-\AscendingNode)$ & $.\ .\ .$ & $9.5348776\,\mathrm n$\\
$\text{Compl.}\ \log\sin(C+\odot-\AscendingNode)$ & $$ & $0.6685194\,\mathrm n$\\
$\text{Compl.}\ \log\cos i$ & $.\ .\ .$ & $0.0075025$\\
\hline
$\log\operatorname{tang}(\upsilon-\AscendingNode)$ & $.\ .\ .$ & $9.3352578\,\mathrm n$\\
\end{tabular}
\end{center}
\[
\begin{array}{llrrr}
\text{also}&C=-&7^{0}&39'&7''056\\
&C+\odot-\AscendingNode=&192&23&15.864
\end{array}
\]
wie vorhin.


\end{minipage}\par
\par\medskip\noindent\begin{minipage}[t]{\textwidth}
\begin{center}$4^{0}$.\end{center}
\begin{center}\small
\begin{tabular}{@{}lcl@{}}
$\log\cos(\odot-\AscendingNode)$ & $.\ .\ .$ & $9.9728762\,\mathrm n$\\
$\log\operatorname{tang}\ObservedBeta$ & $.\ .\ .\ .\ .$ & $8.7349698\,\mathrm n$\\
$\text{Compl.}\ \log\cos(\odot-\alpha)$ & $.$ & $0.1714973\,\mathrm n$\\
$\text{Compl.}\ \log\operatorname{tang}i$ & $.\ .\ .$ & $0.7270128$\\
\hline
$\log\operatorname{tang}D$ & $.\ .\ .\ .\ .$ & $9.6063561\,\mathrm n$\\
\hline
$\log\sin D$ & $.\ .\ .\ .\ .$ & $9.5735295\,\mathrm n$\\
$\log\operatorname{tang}(\odot-\AscendingNode)$ & $.\ .\ .$ & $9.5620014$\\
$\log\cos(\odot-\alpha)$ & $.\ .\ .\ .$ & $9.8285027\,\mathrm n$\\
$\text{Compl.}\ \log\sin(D+\odot-\alpha)$ & $$ & $0.3637217\,\mathrm n$\\
$\text{Compl.}\ \log\cos i$ & $.\ .\ .$ & $0.0075025$\\
\hline
$\log\operatorname{tang}(\upsilon-\AscendingNode)$ & $.\ .\ .\ .$ & $9.3352578\,\mathrm n$\\
\end{tabular}
\end{center}
\[
\begin{array}{llrrr}
\text{also}&D=-&21^{0}&59'&51''182\\
&D+\odot-\alpha=&205&38&41.348
\end{array}
\]
wie oben.


\end{minipage}\par
\par\medskip\noindent\begin{minipage}[t]{\textwidth}
\begin{center}$5^{0}$.\end{center}
\begin{center}\small
\begin{tabular}{@{}lcl@{}}
$\log\operatorname{tang}\ObservedBeta$ & $.\ .\ .\ .\ .$ & $8.7349698\,\mathrm n$\\
$\log\sin(\alpha-\AscendingNode)$ & $.\ .\ .\ .$ & $9.6658973\,\mathrm n$\\
\hline
$\log\operatorname{tang}E$ & $.\ .\ .\ .\ .$ & $9.0690725$\\
\hline
$\log\sin E$ & $.\ .\ .\ .\ .$ & $9.0661081$\\
$\log\sin(\odot-\AscendingNode)$ & $.\ .\ .\ .$ & $9.5348776\,\mathrm n$\\
$\text{Compl.}\ \log\sin(i-E)$ & $.$ & $1.1637907$\\
$\text{Compl.}\ \log\sin(\upsilon-\AscendingNode)$ & $.$ & $0.6746802\,\mathrm n$\\
\hline
$\log\frac{r}{R}$ & $.\ .\ .\ .\ .\ .\ .$ & $0.4394566$\\
\end{tabular}
\end{center}
\[
\begin{array}{llrrr}
\text{also}&E=&6^{0}&41'&12''412\\
&i-E=&3&55&57.138
\end{array}
\]
\begin{center}ferner\end{center}
\[\log r=\log R+\log\frac{r}{R}=0.4320724\]

\end{minipage}\par

\SourceLeaf{96}%

\par\medskip\noindent\begin{minipage}[t]{\textwidth}
\begin{center}$6^{0}$.\end{center}
\begin{center}\small
\begin{tabular}{@{}lcl@{}}
$\log\operatorname{tang}i$ & $.\ .\ .\ .\ .$ & $9.2729872$\\
$\log\sin(\alpha-\AscendingNode)$ & $.\ .\ .\ .$ & $9.6658973\,\mathrm n$\\
\hline
$\log\operatorname{tang}F$ & $.\ .\ .\ .\ .$ & $8.9388845\,\mathrm n$\\
\hline
$\log\cos F$ & $.\ .\ .\ .\ .$ & $9.9983674$\\
$\log\sin\ObservedBeta$ & $.\ .\ .\ .\ .\ .$ & $8.7343300\,\mathrm n$\\
$\log\sin(\odot-\AscendingNode)$ & $.\ .\ .\ .$ & $9.5348776\,\mathrm n$\\
$\text{Compl.}\ \log\sin(F-\ObservedBeta)$ & $.$ & $1.4896990\,\mathrm n$\\
$\text{Compl.}\ \log\sin(\upsilon-\AscendingNode)$ & $.$ & $0.6746802\,\mathrm n$\\
$\text{Compl.}\ \log\cos i$ & $.\ .\ .$ & $0.0075025$\\
\hline
$\log\frac{r}{R}$ & $.\ .\ .\ .\ .\ .\ .$ & $0.4394567$\\
\end{tabular}
\end{center}
\[
\begin{array}{llrrr}
\text{daher}&F=-&4^{0}&57'&53''955\\
&F-\ObservedBeta=-&1&51&20.394
\end{array}
\]
nahe wie vorher.


\end{minipage}\par
\par\medskip\noindent\begin{minipage}[t]{\textwidth}
\begin{center}$7^{0}$.\end{center}
\begin{center}\small
\begin{tabular}{@{}lcl@{}}
$\log\cos i$ & $.\ .\ .\ .\ .\ .$ & $9.9924975$\\
$\log\operatorname{tang}(\upsilon-\AscendingNode)$ & $.\ .\ .$ & $9.3352577\,\mathrm n$\\
\hline
$\log\operatorname{tang}G$ & $.\ .\ .\ .\ .$ & $9.3277552\,\mathrm n$\\
\hline
$\log\cos G$ & $.\ .\ ,\ .\ .$ & $9.9903922$\\
$\log\sin(\odot-\alpha)$ & $.\ .\ .\ .$ & $9.8686173\,\mathrm n$\\
$\text{Compl.}\ \log\sin(\alpha-\AscendingNode-G)$ & $$ & $0.5705092\,\mathrm n$\\
$\text{Compl.}\ \log\cos(\upsilon-\AscendingNode)$ & $.$ & $0.0099379$\\
\hline
$\log\frac{r}{R}$ & $.\ .\ .\ .\ .\ .\ .$ & $0.4394566$\\
\end{tabular}
\end{center}
\[
\begin{array}{llrrr}
\text{also}&G=-&12^{0}&0'&27''118\\
&\alpha-\AscendingNode-G=-&15&35&42.492
\end{array}
\]
wie oben.


\end{minipage}\par
\par\medskip\noindent\begin{minipage}[t]{\textwidth}
\begin{center}$8^{0}$.\end{center}
\begin{center}\small
\begin{tabular}{@{}lcl@{}}
$\log\operatorname{tang}(\alpha-\AscendingNode)$ & $.\ .\ .$ & $9.7183744\,\mathrm n$\\
$\log\cos i$ & $.\ .\ .\ .\ .\ .$ & $9.9924975$\\
\hline
$\log\operatorname{tang}H$ & $.\ .\ .\ .\ .$ & $9.7258769\,\mathrm n$\\
\hline
$\log\sin H$ & $.\ .\ .\ .\ .$ & $9.6717672\,\mathrm n$\\
$\log\sin(\odot-\alpha)$ & $.\ .\ .\ .$ & $9.8686173\,\mathrm n$\\
$\text{Compl.}\ \log\sin(H-(\upsilon-\AscendingNode))$ & $$ & $0.5649695\,\mathrm n$\\
$\text{Compl.}\ \log\sin(\alpha-\AscendingNode)$ & $.$ & $0.3341027\,\mathrm n$\\
\hline
$\log\frac{r}{R}$ & $.\ .\ .\ .\ .\ .\ .$ & $0.4394567$\\
\end{tabular}
\end{center}
\[
\begin{array}{llrrr}
\text{folglich}&H=-&28^{0}&0'&39''879\\
&H-(\upsilon-\AscendingNode)=-&15&48&1.937
\end{array}
\]
wie vorher.

\end{minipage}\par

\SourceLeaf{97}%

\par\medskip\noindent\begin{minipage}[t]{\textwidth}
\begin{center}$9^{0}$.\end{center}
\begin{center}\small
\begin{tabular}{@{}lcl@{}}
$\log\operatorname{tang}\ObservedBeta$ & $.\ .\ .\ .\ .$ & $8.7349698\,\mathrm n$\\
$\text{Compl.}\ \log\sin i$ & $.\ .\ .$ & $0.7345153$\\
$\text{Compl.}\ \log\cos(\alpha-\AscendingNode)$ & $.$ & $0.0524771$\\
\hline
$\log\operatorname{tang}I$ & $.\ .\ .\ .\ .$ & $9.5219622\,\mathrm n$\\
\hline
$\log\sin I$ & $.\ .\ .\ .\ .\ .$ & $9.4991749\,\mathrm n$\\
$\log\cos(\odot-\AscendingNode)$ & $.\ .\ .$ & $9.9728762\,\mathrm n$\\
$\text{Compl.}\ \log\sin(\upsilon-\AscendingNode-I)$ & $$ & $0.9674054$\\
\hline
$\log\frac{r}{R}$ & $.\ .\ .\ .\ .\ .\ .$ & $0.4394565$\\
\end{tabular}
\end{center}
\[
\begin{array}{llrrr}
\text{hieraus}&I=-&18^{0}&23'&55''334\\
&\upsilon-\AscendingNode-I=&6&11&17.392
\end{array}
\]
wie vorhin.


\end{minipage}\par
\par\medskip\noindent\begin{minipage}[t]{\textwidth}
\begin{center}$10^{0}$.\end{center}
\begin{center}In der Nähe des Knotens weniger scharf.\end{center}
\begin{center}\small
\begin{tabular}{@{}lcl@{}}
$\log\sin i$ & $.\ .\ .\ .\ .\ .$ & $9.2654847$\\
$\log\cos(\alpha-\AscendingNode)$ & $.\ .\ .\ .$ & $9.9475229$\\
$\log\operatorname{tang}(\upsilon-\AscendingNode)$ & $.\ .\ .\ .$ & $9.3352577\,\mathrm n$\\
\hline
$\log\operatorname{tang}K$ & $.\ .\ .\ .\ .$ & $8.5482653\,\mathrm n$\\
\hline
$\log\cos K$ & $.\ .\ .\ .\ .$ & $9.9997290$\\
$\log\sin\ObservedBeta$ & $.\ .\ .\ .\ .\ .$ & $8.7343300\,\mathrm n$\\
$\log\cos(\odot-\AscendingNode)$ & $.\ .\ .$ & $9.9728762\,\mathrm n$\\
$\text{Compl.}\ \log\sin(K-\ObservedBeta)$ & $.$ & $1.7225836$\\
$\text{Compl.}\ \log\cos(\upsilon-\AscendingNode)$ & $.$ & $0.0099379$\\
\hline
$\log\frac{r}{R}$ & $.\ .\ .\ .\ .\ .\ .$ & $0.4394567$\\
\end{tabular}
\end{center}
\[
\begin{array}{llrrr}
\text{also}&K=-&2^{0}&1'&26''344\\
&K-\ObservedBeta=&1&5&7.217
\end{array}
\]
wie vorhin.


\end{minipage}\par
\par\medskip\noindent\begin{minipage}[t]{\textwidth}
\begin{center}$11^{0}$.\end{center}
\[C+\odot-\alpha=219^{0}\ 59'\ 25''474\]
\begin{center}\small
\begin{tabular}{@{}lcl@{}}
$\log\sin C$ & $.\ .\ .\ .\ .$ & $9.1243583\,\mathrm n$\\
$\log\sin(\odot-\alpha)$ & $.\ .\ .\ .$ & $9.8686173\,\mathrm n$\\
$\text{Compl.}\ \log\cos(C+\odot-\alpha)$ & $$ & $0.1156850\,\mathrm n$\\
$\text{Compl.}\ \log\operatorname{tang}(\odot-\AscendingNode)$ & $$ & $0.4379986$\\
$\text{Compl.}\ \log\cos i$ & $.\ .\ .$ & $0.0075025$\\
\hline
$\log\operatorname{tang}L$ & $.\ .\ .\ .\ .$ & $9.5541617\,\mathrm n$\\
\hline
$\log\sin L$ & $.\ .\ .\ .\ .\ .$ & $9.5279439\,\mathrm n$\\
$\text{Compl.}\ \log\sin(\upsilon-\AscendingNode-L)$ & $$ & $0.8843888$\\
$\text{Compl.}\ \log\cos(\odot-\AscendingNode)$ & $.$ & $0.0271238\,\mathrm n$\\
\hline
$\log\frac{r}{R}$ & $.\ .\ .\ .\ .\ .\ .$ & $0.4394565$\\
\end{tabular}
\end{center}
\[
\begin{array}{llrrr}
\text{also}&L=-&19^{0}&42'&32''533\\
&\upsilon-\AscendingNode-L=&7&29&54.591
\end{array}
\]
wie zuvor.

\end{minipage}\par

\SourceLeaf{98}%

\par\medskip\noindent\begin{minipage}[t]{\textwidth}
\begin{center}$12^{0}$.\end{center}
\[D+\odot-\AscendingNode=178^{0}\ 2'\ 31''738\]
\begin{center}\small
\begin{tabular}{@{}lcl@{}}
$\log\sin D$ & $.\ .\ .\ .\ .$ & $9.5735295\,\mathrm n$\\
$\log\cos(\odot-\AscendingNode)$ & $.\ .\ .$ & $9.9728762\,\mathrm n$\\
$\text{Compl.}\ \log\cos(D+\odot-\AscendingNode)$ & $$ & $0.0002536\,\mathrm n$\\
$\text{Compl.}\ \log\cos i$ & $.\ .\ .$ & $0.0075025$\\
\hline
$\log\operatorname{tang}(M=L)$ & $.\ .\ .$ & $9.5541618\,\mathrm n$\\
\end{tabular}
\end{center}
wie oben in $11^{0}$.

Der übrige Theil der Rechnung eben so wie dort.


\end{minipage}\par
\par\medskip\noindent\begin{minipage}[t]{\textwidth}
\begin{center}$13^{0}$.\end{center}
\begin{center}\small
\begin{tabular}{@{}lcl@{}}
$\log r$ & $.\ .\ .\ .\ .\ .\ .$ & $0.4320724$\\
$\log\sin(\upsilon-\AscendingNode)$ & $.\ .\ .\ .$ & $9.3253198\,\mathrm n$\\
$\log\sin i$ & $.\ .\ .\ .\ .\ .$ & $9.2654847$\\
$\text{Compl.}\ \log\sin\ObservedBeta$ & $.\ .\ .$ & $1.2656700\,\mathrm n$\\
\hline
$\log\Delta$ & $.\ .\ .\ .\ .\ .\ .$ & $0.2885469$\\
\end{tabular}
\end{center}
\par\medskip\noindent\rule{.3\textwidth}{.3pt}

\end{minipage}\par
